\documentclass[11pt,a4paper,openany]{book}

%===========================================================
% IDIOMA Y CODIFICACIÓN
%===========================================================
\usepackage[spanish,es-noquoting]{babel}
\usepackage[utf8]{inputenc}
\usepackage[T1]{fontenc}

%===========================================================
% TIPOGRAFÍA
%===========================================================
\usepackage{lmodern}
\usepackage{microtype}

%===========================================================
% GEOMETRÍA DE PÁGINA
%===========================================================
\usepackage{geometry}
\geometry{
    top=2.7cm,
    bottom=2.7cm,
    left=3cm,
    right=3cm
}

%===========================================================
% MATEMÁTICA
%===========================================================
\usepackage{amsmath, amssymb, amsthm, mathtools}
\usepackage{mathrsfs}

%===========================================================
% GRÁFICOS
%===========================================================
\usepackage{tikz}
\usetikzlibrary{
    arrows.meta,
    positioning,
    calc,
    patterns,
    decorations.pathreplacing
}

%===========================================================
% COLORES
%===========================================================
\usepackage{xcolor}

\definecolor{deepblue}{RGB}{18,45,105}
\definecolor{softyellow}{RGB}{248,236,170}
\definecolor{theoremgray}{RGB}{235,235,235}
\definecolor{linegray}{RGB}{150,150,150}

%===========================================================
% HIPERVÍNCULOS
%===========================================================
\usepackage{hyperref}
\hypersetup{
    colorlinks=true,
    linkcolor=deepblue,
    citecolor=deepblue,
    urlcolor=deepblue
}

%===========================================================
% LISTAS
%===========================================================
\usepackage{enumitem}
\setlist[itemize]{topsep=4pt,itemsep=2pt}
\setlist[enumerate]{topsep=4pt,itemsep=2pt}

%===========================================================
% CAJAS TIPO LIBRO
%===========================================================
\usepackage[most]{tcolorbox}

\newtcolorbox{grisplano}{
  colback=gray!18,
  colframe=gray!18,
  boxrule=0pt,
  arc=0mm,
  left=10pt,
  right=10pt,
  top=8pt,
  bottom=8pt,
  before skip=10pt,
  after skip=10pt
}

\newtcolorbox{bloquelibro}{
  colback=gray!12,
  colframe=gray!45,
  boxrule=0.4pt,
  arc=0mm,
  left=10pt,
  right=10pt,
  top=8pt,
  bottom=8pt,
  before skip=10pt,
  after skip=10pt
}

\newtcolorbox{comentariolibro}{
  colback=softyellow!45,
  colframe=softyellow!80!black,
  boxrule=0.4pt,
  arc=0mm,
  left=10pt,
  right=10pt,
  top=7pt,
  bottom=7pt,
  before skip=10pt,
  after skip=10pt
}

%===========================================================
% ENTORNOS MATEMÁTICOS
%===========================================================
\theoremstyle{definition}
\newtheorem{definicion}{Definición}[chapter]
\newtheorem{ejemplo}{Ejemplo}[chapter]

\theoremstyle{plain}
\newtheorem{teorema}{Teorema}[chapter]
\newtheorem{proposicion}{Proposición}[chapter]
\newtheorem{lema}{Lema}[chapter]
\newtheorem{corolario}{Corolario}[chapter]

\theoremstyle{remark}
\newtheorem*{observacionthm}{Observación}

%===========================================================
% ENTORNOS PERSONALIZADOS
%===========================================================
\renewenvironment{proof}[1][\textbf{Demostración.}]{%
  \par\noindent #1\ \normalfont
}{\hfill$\square$\par\medskip}

\newenvironment{observacion}
{\par\medskip\noindent\textbf{Observación. }\normalfont}
{\par\medskip}

\newenvironment{notacion}
{\par\medskip\noindent\textbf{Notación. }\normalfont}
{\par\medskip}

\newenvironment{idea}
{\begin{comentariolibro}\textbf{Idea de la demostración. }\normalfont}
{\end{comentariolibro}}

%===========================================================
% COMANDOS ÚTILES
%===========================================================
\newcommand{\R}{\mathbb{R}}
\newcommand{\Q}{\mathbb{Q}}
\newcommand{\N}{\mathbb{N}}
\newcommand{\Z}{\mathbb{Z}}
\newcommand{\C}{\mathbb{C}}

\newcommand{\eps}{\varepsilon}
\newcommand{\ds}{\displaystyle}

% Ínfimo y supremo en español
\DeclareMathOperator{\sen}{sen}
\DeclareMathOperator{\arcsen}{arcsen}
\DeclareMathOperator{\Dom}{Dom}

% Integral inferior y superior
\newcommand{\intinf}[2]{\underline{\int_{#1}^{#2}}}
\newcommand{\intsup}[2]{\overline{\int_{#1}^{#2}}}

%===========================================================
% ESTILO DE CAPÍTULOS Y SECCIONES
%===========================================================
\usepackage{titlesec}

\titleformat{\chapter}[display]
  {\normalfont\huge\bfseries\color{deepblue}}
  {\chaptername\ \thechapter}
  {10pt}
  {\Huge}

\titleformat{\section}
  {\normalfont\Large\bfseries\color{deepblue}}
  {\thesection}
  {1em}
  {}

\titleformat{\subsection}
  {\normalfont\large\bfseries\color{deepblue}}
  {\thesubsection}
  {1em}
  {}

%===========================================================
% DATOS DEL DOCUMENTO
%===========================================================
\title{Análisis Real\\\large Sucesiones y series de funciones reales}
\author{Miguel Angel Tintaya Condori}
\date{2026}

%===========================================================
% INICIO DEL DOCUMENTO
%===========================================================
\begin{document}

\maketitle
\tableofcontents
%===========================================================
\chapter{Sucesiones y series de funciones}
%===========================================================

\section{Convergencia puntual y convergencia uniforme}

En capítulos anteriores estudiamos sucesiones de números reales. Ahora
estudiaremos sucesiones cuyos términos ya no son números, sino funciones.

La idea general es la siguiente: muchas veces queremos construir una función
$f$ con ciertas propiedades, pero lo que obtenemos primero es una sucesión de
funciones
\[
    f_1,\ f_2,\ \ldots,\ f_n,\ \ldots
\]
que se acercan cada vez más a una función límite. El problema central será
entender en qué sentido se acercan.

\begin{observacion}
Para sucesiones numéricas existe una sola noción usual de convergencia.
Para sucesiones de funciones existen varias. Las dos más importantes para
nuestro curso serán:

\begin{itemize}
    \item convergencia puntual;
    \item convergencia uniforme.
\end{itemize}

La convergencia uniforme será más fuerte y será la que conserve propiedades
como continuidad, integrabilidad y derivabilidad bajo ciertas hipótesis.
\end{observacion}

%-----------------------------------------------------------
\subsection{Convergencia puntual}
%-----------------------------------------------------------

\begin{grisplano}
\begin{definicion}[Convergencia puntual]
Sea $X\subset \R$ y sea
\[
    (f_n)_{n\in\N},\qquad f_n:X\to\R,
\]
una sucesión de funciones. Decimos que $(f_n)$ \emph{converge puntualmente}
a una función $f:X\to\R$ si, para cada $x\in X$, la sucesión numérica
\[
    \bigl(f_n(x)\bigr)_{n\in\N}
\]
converge a $f(x)$.

Es decir,
\[
    f_n\to f \text{ puntualmente en }X
\]
si y solo si
\[
    \forall x\in X,\qquad
    \lim_{n\to\infty} f_n(x)=f(x).
\]
\end{definicion}
\end{grisplano}

En términos de $\varepsilon$, esto significa que
\[
    \forall x\in X,\ \forall \varepsilon>0,\ \exists n_0\in\N
\]
tal que
\[
    n\geq n_0
    \quad\Longrightarrow\quad
    |f_n(x)-f(x)|<\varepsilon.
\]

\begin{observacion}
En la convergencia puntual, el entero $n_0$ puede depender de dos cosas:

\[
    n_0=n_0(\varepsilon,x).
\]

Es decir, para cada punto $x$ podemos necesitar un tiempo distinto para que
$f_n(x)$ se acerque a $f(x)$.
\end{observacion}

\begin{observacion}
Geométricamente, fijar $x\in X$ significa mirar una recta vertical. Sobre esa
recta tenemos los puntos
\[
    (x,f_1(x)),\ (x,f_2(x)),\ldots,\ (x,f_n(x)),\ldots
\]
y la convergencia puntual dice que estos puntos se acercan a
\[
    (x,f(x)).
\]
\end{observacion}

\begin{center}
\begin{tikzpicture}[scale=1.1]
    % axes
    \draw[->] (-0.2,0) -- (5.2,0) node[right] {$x$};
    \draw[->] (0,-0.2) -- (0,3.2) node[above] {$y$};

    % vertical line
    \draw[dashed,gray] (2.7,0) -- (2.7,3);

    % points
    \fill[blue] (2.7,2.5) circle (1.7pt) node[right] {$f_1(x)$};
    \fill[blue] (2.7,2.1) circle (1.7pt) node[right] {$f_2(x)$};
    \fill[blue] (2.7,1.75) circle (1.7pt) node[right] {$f_3(x)$};
    \fill[blue] (2.7,1.45) circle (1.7pt) node[right] {$f_n(x)$};

    \fill[red] (2.7,1.1) circle (2pt) node[right] {$f(x)$};

    \node[below] at (2.7,0) {$x$};
\end{tikzpicture}
\end{center}

\begin{ejemplo}
Sea
\[
    f_n:\R\to\R,\qquad f_n(x)=\frac{x}{n}.
\]
Entonces
\[
    f_n\to 0
\]
puntualmente en $\R$.

En efecto, fijado $x\in\R$, tenemos
\[
    \lim_{n\to\infty} f_n(x)
    =
    \lim_{n\to\infty}\frac{x}{n}
    =
    0.
\]
Por tanto, la función límite es
\[
    f(x)=0,\qquad x\in\R.
\]
\end{ejemplo}

%-----------------------------------------------------------
\subsection{Convergencia uniforme}
%-----------------------------------------------------------

\begin{grisplano}
\begin{definicion}[Convergencia uniforme]
Sea $X\subset\R$ y sea $(f_n)_{n\in\N}$ una sucesión de funciones
\[
    f_n:X\to\R.
\]
Decimos que $(f_n)$ \emph{converge uniformemente} a $f:X\to\R$ en $X$ si
\[
    \forall \varepsilon>0,\ \exists n_0\in\N
\]
tal que
\[
    n\geq n_0
    \quad\Longrightarrow\quad
    |f_n(x)-f(x)|<\varepsilon
    \quad\text{para todo }x\in X.
\]
\end{definicion}
\end{grisplano}

\begin{observacion}
La diferencia con la convergencia puntual está en el orden de los cuantificadores.

En la convergencia puntual:
\[
    \forall x\in X,\ \forall\varepsilon>0,\ \exists n_0=n_0(\varepsilon,x).
\]

En la convergencia uniforme:
\[
    \forall\varepsilon>0,\ \exists n_0=n_0(\varepsilon),\ \forall x\in X.
\]

En la convergencia uniforme, el mismo $n_0$ sirve para todos los puntos de $X$.
\end{observacion}

\begin{observacion}
Geométricamente, dado $\varepsilon>0$, consideramos la banda alrededor del
gráfico de $f$:
\[
    B_\varepsilon(f)
    =
    \{(x,y)\in X\times\R:
    f(x)-\varepsilon<y<f(x)+\varepsilon\}.
\]
Decir que $f_n\to f$ uniformemente significa que, a partir de cierto $n_0$,
todos los gráficos de $f_n$ quedan completamente dentro de esa banda.
\end{observacion}

\begin{center}
\begin{tikzpicture}[scale=1.05]
    \draw[->] (-0.2,0) -- (6.2,0) node[right] {$x$};
    \draw[->] (0,-0.2) -- (0,3.7) node[above] {$y$};

    % band
    \fill[gray!20,domain=0.5:5.8,smooth,variable=\x]
        plot ({\x},{1.8+0.35*sin(90*\x)+0.45})
        -- plot[domain=5.8:0.5] ({\x},{1.8+0.35*sin(90*\x)-0.45})
        -- cycle;

    % f
    \draw[blue,thick,domain=0.5:5.8,smooth,variable=\x]
        plot ({\x},{1.8+0.35*sin(90*\x)})
        node[right] {$f$};

    % epsilon lines
    \draw[gray,dashed,domain=0.5:5.8,smooth,variable=\x]
        plot ({\x},{1.8+0.35*sin(90*\x)+0.45});
    \draw[gray,dashed,domain=0.5:5.8,smooth,variable=\x]
        plot ({\x},{1.8+0.35*sin(90*\x)-0.45});

    % fn
    \draw[red,thick,domain=0.5:5.8,smooth,variable=\x]
        plot ({\x},{1.8+0.35*sin(90*\x)+0.2*cos(150*\x)})
        node[right] {$f_n$};

    \node at (4.9,3.1) {$+\varepsilon$};
    \node at (4.9,0.8) {$-\varepsilon$};
\end{tikzpicture}
\end{center}

%-----------------------------------------------------------
\subsection{Primeros ejemplos}
%-----------------------------------------------------------

\begin{ejemplo}
La sucesión
\[
    f_n:\R\to\R,\qquad f_n(x)=\frac{x}{n},
\]
converge puntualmente a la función nula, pero no converge uniformemente en
$\R$.

\begin{proof}
Ya sabemos que, para cada $x\in\R$,
\[
    \lim_{n\to\infty}\frac{x}{n}=0.
\]
Luego $f_n\to 0$ puntualmente en $\R$.

Veamos que la convergencia no es uniforme. Si fuera uniforme, para
$\varepsilon=1$ existiría $n_0\in\N$ tal que
\[
    n\geq n_0
    \quad\Longrightarrow\quad
    \left|\frac{x}{n}\right|<1
    \quad\text{para todo }x\in\R.
\]
Tomando $n=n_0$ y $x=2n_0$, obtenemos
\[
    \left|\frac{x}{n_0}\right|
    =
    \left|\frac{2n_0}{n_0}\right|
    =
    2,
\]
lo cual contradice la desigualdad anterior. Por tanto, la convergencia no es
uniforme en $\R$.
\end{proof}
\end{ejemplo}

\begin{observacion}
El problema anterior aparece porque $\R$ no es acotado. En cambio, si
$X\subset\R$ es acotado y existe $c>0$ tal que
\[
    |x|\leq c
    \qquad \text{para todo }x\in X,
\]
entonces
\[
    \frac{x}{n}\to 0
\]
uniformemente en $X$.

En efecto, dado $\varepsilon>0$, tomamos $n_0\in\N$ tal que
\[
    n_0>\frac{c}{\varepsilon}.
\]
Entonces, si $n\geq n_0$ y $x\in X$,
\[
    \left|\frac{x}{n}\right|
    \leq \frac{c}{n}
    \leq \frac{c}{n_0}
    <\varepsilon.
\]
\end{observacion}

\begin{ejemplo}
Sea
\[
    f_n:[0,1]\to\R,\qquad f_n(x)=x^n.
\]
Entonces $(f_n)$ converge puntualmente a la función
\[
    f(x)=
    \begin{cases}
        0, & 0\leq x<1,\\
        1, & x=1.
    \end{cases}
\]
La convergencia no es uniforme en $[0,1]$.

\begin{proof}
Si $0\leq x<1$, entonces
\[
    \lim_{n\to\infty}x^n=0.
\]
Si $x=1$, entonces
\[
    \lim_{n\to\infty}1^n=1.
\]
Por tanto, la función límite puntual es
\[
    f(x)=
    \begin{cases}
        0, & 0\leq x<1,\\
        1, & x=1.
    \end{cases}
\]

Ahora probemos que la convergencia no es uniforme. Tomemos
\[
    \varepsilon=\frac12.
\]
Si la convergencia fuese uniforme, existiría $n_0\in\N$ tal que, para todo
$n\geq n_0$ y todo $x\in[0,1]$,
\[
    |x^n-f(x)|<\frac12.
\]
Fijemos $n=n_0$. Como
\[
    \lim_{x\to 1^-}x^{n_0}=1,
\]
existe $x\in[0,1)$ suficientemente cercano a $1$ tal que
\[
    x^{n_0}>\frac12.
\]
Pero para ese $x$ se tiene $f(x)=0$, y entonces
\[
    |f_{n_0}(x)-f(x)|
    =
    |x^{n_0}-0|
    =
    x^{n_0}
    >
    \frac12.
\]
Esto contradice la convergencia uniforme.
\end{proof}
\end{ejemplo}

\begin{center}
\begin{tikzpicture}[scale=4]
    \draw[->] (-0.05,0) -- (1.12,0) node[right] {$x$};
    \draw[->] (0,-0.05) -- (0,1.15) node[above] {$y$};

    \draw[domain=0:1,smooth,variable=\x,blue,thick] plot ({\x},{\x});
    \draw[domain=0:1,smooth,variable=\x,red,thick] plot ({\x},{\x*\x});
    \draw[domain=0:1,smooth,variable=\x,green!60!black,thick] plot ({\x},{\x*\x*\x});
    \draw[domain=0:1,smooth,variable=\x,purple,thick] plot ({\x},{\x*\x*\x*\x});

    \node[blue] at (0.78,0.82) {$x$};
    \node[red] at (0.80,0.58) {$x^2$};
    \node[green!60!black] at (0.82,0.43) {$x^3$};
    \node[purple] at (0.83,0.30) {$x^4$};

    \node[below] at (1,0) {$1$};
    \node[left] at (0,1) {$1$};
\end{tikzpicture}
\end{center}

\begin{observacion}
Aunque $f_n(x)=x^n$ no converge uniformemente en $[0,1]$, sí converge
uniformemente en todo intervalo
\[
    [0,1-\delta],
    \qquad 0<\delta<1.
\]
En efecto, si $0\leq x\leq 1-\delta$, entonces
\[
    0\leq x^n\leq (1-\delta)^n,
\]
y como $0<1-\delta<1$, se tiene
\[
    (1-\delta)^n\to 0.
\]
Así,
\[
    \sup_{x\in[0,1-\delta]}|x^n-0|
    \leq (1-\delta)^n\to 0.
\]
\end{observacion}

\begin{ejemplo}
Sea
\[
    f_n:[0,1]\to\R,\qquad f_n(x)=x^n(1-x^n).
\]
Entonces $f_n\to 0$ puntualmente en $[0,1]$, pero la convergencia no es
uniforme.

\begin{proof}
Si $0\leq x<1$, entonces $x^n\to 0$, de donde
\[
    x^n(1-x^n)\to 0.
\]
Si $x=1$, entonces
\[
    f_n(1)=1^n(1-1^n)=0.
\]
Por tanto, $f_n\to 0$ puntualmente en $[0,1]$.

Veamos que la convergencia no es uniforme. Para cada $n$, tomamos
\[
    x_n=\sqrt[n]{\frac12}.
\]
Entonces $x_n\in[0,1]$ y
\[
    x_n^n=\frac12.
\]
Por tanto,
\[
    f_n(x_n)
    =
    x_n^n(1-x_n^n)
    =
    \frac12\left(1-\frac12\right)
    =
    \frac14.
\]
Así,
\[
    \sup_{x\in[0,1]}|f_n(x)-0|
    \geq
    |f_n(x_n)|
    =
    \frac14.
\]
Luego la distancia uniforme entre $f_n$ y $0$ no tiende a cero. Por tanto,
$f_n$ no converge uniformemente a $0$.
\end{proof}
\end{ejemplo}

%-----------------------------------------------------------
\section{Propiedades de la convergencia uniforme}
%-----------------------------------------------------------

\subsection{La convergencia uniforme conserva la continuidad}

\begin{grisplano}
\begin{teorema}[Límite uniforme de funciones continuas]
Sea $(f_n)$ una sucesión de funciones
\[
    f_n:X\to\R.
\]
Supongamos que
\[
    f_n\to f
\]
uniformemente en $X$. Si cada $f_n$ es continua en un punto $a\in X$, entonces
$f$ es continua en $a$.
\end{teorema}
\end{grisplano}

\begin{idea}
La prueba usa el truco clásico de sumar y restar $f_n(x)$ y $f_n(a)$:
\[
    f(x)-f(a)
    =
    [f(x)-f_n(x)]
    +
    [f_n(x)-f_n(a)]
    +
    [f_n(a)-f(a)].
\]
Los términos de los extremos se controlan por convergencia uniforme. El término
del centro se controla por continuidad de $f_n$.
\end{idea}

\begin{proof}
Sea $\varepsilon>0$. Como $f_n\to f$ uniformemente en $X$, existe
$n_0\in\N$ tal que
\[
    n\geq n_0
    \quad\Longrightarrow\quad
    |f_n(x)-f(x)|<\frac{\varepsilon}{3}
    \quad\text{para todo }x\in X.
\]
Fijemos un índice $n\geq n_0$. Como $f_n$ es continua en $a$, existe
$\delta>0$ tal que
\[
    x\in X,\ |x-a|<\delta
    \quad\Longrightarrow\quad
    |f_n(x)-f_n(a)|<\frac{\varepsilon}{3}.
\]
Entonces, si $x\in X$ y $|x-a|<\delta$, tenemos
\[
\begin{aligned}
    |f(x)-f(a)|
    &\leq |f(x)-f_n(x)|
        +|f_n(x)-f_n(a)|
        +|f_n(a)-f(a)|  \\
    &< \frac{\varepsilon}{3}
        +\frac{\varepsilon}{3}
        +\frac{\varepsilon}{3}
    =\varepsilon.
\end{aligned}
\]
Luego $f$ es continua en $a$.
\end{proof}

\begin{corolario}
Si $f_n:X\to\R$ son continuas en todo $X$ y $f_n\to f$ uniformemente en $X$,
entonces $f$ es continua en $X$.
\end{corolario}

\begin{observacion}
Este resultado explica por qué la convergencia de
\[
    f_n(x)=x^n
\]
a
\[
    f(x)=
    \begin{cases}
        0, & 0\leq x<1,\\
        1, & x=1,
    \end{cases}
\]
no puede ser uniforme en $[0,1]$: cada $f_n$ es continua, pero el límite puntual
$f$ no es continuo.
\end{observacion}

%-----------------------------------------------------------
\subsection{Convergencia monótona de funciones}
%-----------------------------------------------------------

\begin{grisplano}
\begin{definicion}[Convergencia monótona]
Sea $X\subset\R$ y sea $(f_n)$ una sucesión de funciones
\[
    f_n:X\to\R.
\]
Decimos que $f_n\to f$ \emph{monótonamente} si, para cada $x\in X$, la sucesión
numérica
\[
    (f_n(x))_{n\in\N}
\]
es monótona y converge a $f(x)$.
\end{definicion}
\end{grisplano}

\begin{observacion}
Por ejemplo, si para todo $x\in X$ se tiene
\[
    f_1(x)\leq f_2(x)\leq \cdots \leq f_n(x)\leq \cdots
\]
y
\[
    f_n(x)\to f(x),
\]
entonces
\[
    f_n\to f
\]
monótonamente.

También puede ocurrir el caso decreciente:
\[
    f_1(x)\geq f_2(x)\geq \cdots \geq f_n(x)\geq \cdots.
\]
\end{observacion}

\begin{observacion}
Si $f_n\to f$ monótonamente, entonces para todo $x\in X$,
\[
    |f_{n+1}(x)-f(x)|
    \leq
    |f_n(x)-f(x)|.
\]
Es decir, el error puntual disminuye cuando aumenta $n$.
\end{observacion}

%-----------------------------------------------------------
\subsection{Teorema de Dini}
%-----------------------------------------------------------

\begin{grisplano}
\begin{teorema}[Dini]
Sea $X\subset\R$ compacto. Sea
\[
    (f_n)_{n\in\N},\qquad f_n:X\to\R,
\]
una sucesión de funciones continuas. Supongamos que
\[
    f_n\to f
\]
monótonamente en $X$ y que $f:X\to\R$ es continua. Entonces
\[
    f_n\to f
\]
uniformemente en $X$.
\end{teorema}
\end{grisplano}

\begin{idea}
Queremos probar que, dado $\varepsilon>0$, a partir de cierto índice ya no
existe ningún punto $x\in X$ donde el error sea mayor o igual que
$\varepsilon$.

Para eso definimos los conjuntos malos:
\[
    K_n=\{x\in X:|f_n(x)-f(x)|\geq \varepsilon\}.
\]
Estos conjuntos son compactos, están encajados y su intersección es vacía.
Por compacidad, alguno debe ser vacío.
\end{idea}

\begin{proof}
Sea $\varepsilon>0$. Para cada $n\in\N$, definimos
\[
    K_n
    =
    \{x\in X: |f_n(x)-f(x)|\geq \varepsilon\}.
\]

Como $f_n$ y $f$ son continuas, la función
\[
    x\mapsto |f_n(x)-f(x)|
\]
es continua. Entonces $K_n$ es cerrado en $X$. Como $X$ es compacto, concluimos
que cada $K_n$ es compacto.

Además, por la convergencia monótona, el error
\[
    |f_n(x)-f(x)|
\]
disminuye al crecer $n$. Por tanto,
\[
    K_1\supset K_2\supset K_3\supset \cdots.
\]

Ahora veamos que
\[
    \bigcap_{n=1}^{\infty}K_n=\varnothing.
\]
En efecto, si existiera $p\in\bigcap_{n=1}^{\infty}K_n$, entonces
\[
    |f_n(p)-f(p)|\geq \varepsilon
    \qquad \text{para todo }n\in\N.
\]
Pero esto contradice que
\[
    f_n(p)\to f(p).
\]
Por tanto,
\[
    \bigcap_{n=1}^{\infty}K_n=\varnothing.
\]

Como $(K_n)$ es una sucesión decreciente de compactos no vacíos solo podría
tener intersección no vacía. Luego, al tener intersección vacía, necesariamente
existe $n_0\in\N$ tal que
\[
    K_{n_0}=\varnothing.
\]
Como la sucesión es decreciente,
\[
    K_n=\varnothing
    \qquad \text{para todo }n\geq n_0.
\]
Esto significa que, si $n\geq n_0$, entonces no existe $x\in X$ tal que
\[
    |f_n(x)-f(x)|\geq \varepsilon.
\]
Por tanto,
\[
    |f_n(x)-f(x)|<\varepsilon
    \qquad \text{para todo }x\in X.
\]
Así, $f_n\to f$ uniformemente en $X$.
\end{proof}

\begin{observacion}
La compacidad de $X$ es esencial. Por ejemplo,
\[
    f_n:[0,1)\to\R,\qquad f_n(x)=x^n,
\]
converge monótonamente a $0$ en $[0,1)$, y el límite es continuo. Pero la
convergencia no es uniforme porque $[0,1)$ no es compacto.
\end{observacion}

%-----------------------------------------------------------
\section{Paso al límite bajo el signo integral}
%-----------------------------------------------------------

\begin{grisplano}
\begin{teorema}[Paso al límite bajo el signo integral]
Sea $(f_n)$ una sucesión de funciones integrables
\[
    f_n:[a,b]\to\R.
\]
Si
\[
    f_n\to f
\]
uniformemente en $[a,b]$, entonces $f$ es integrable y
\[
    \int_a^b f(x)\,dx
    =
    \lim_{n\to\infty}\int_a^b f_n(x)\,dx.
\]
Equivalentemente,
\[
    \lim_{n\to\infty}\int_a^b f_n(x)\,dx
    =
    \int_a^b\left(\lim_{n\to\infty}f_n(x)\right)\,dx.
\]
\end{teorema}
\end{grisplano}

\begin{idea}
La prueba tiene dos partes.

Primero probamos que $f$ es integrable. Como $f_n\to f$ uniformemente, para
un índice fijo grande $m$, la función $f$ está muy cerca de $f_m$. Como $f_m$
ya es integrable, existe una partición donde la suma de oscilaciones de $f_m$
es pequeña. La cercanía uniforme permite transferir esa pequeña oscilación de
$f_m$ hacia $f$.

Luego probamos la igualdad de integrales usando
\[
    \left|\int_a^b f-\int_a^b f_n\right|
    =
    \left|\int_a^b(f-f_n)\right|
    \leq
    \int_a^b|f-f_n|.
\]
\end{idea}

\begin{proof}
Sea $\varepsilon>0$. Como $f_n\to f$ uniformemente en $[a,b]$, existe
$n_0\in\N$ tal que
\[
    n\geq n_0
    \quad\Longrightarrow\quad
    |f_n(x)-f(x)|<\frac{\varepsilon}{4(b-a)}
    \qquad\text{para todo }x\in[a,b].
\]
Fijemos un índice
\[
    m\geq n_0.
\]
Entonces
\[
    |f_m(x)-f(x)|<\frac{\varepsilon}{4(b-a)}
    \qquad\text{para todo }x\in[a,b].
\]

Como $f_m$ es integrable, existe una partición
\[
    P=\{a=t_0<t_1<\cdots<t_r=b\}
\]
tal que
\[
    \sum_{i=1}^r \omega_i(f_m)(t_i-t_{i-1})
    <
    \frac{\varepsilon}{2},
\]
donde
\[
    \omega_i(f_m)
    =
    \sup_{x,y\in[t_{i-1},t_i]}
    |f_m(x)-f_m(y)|
\]
es la oscilación de $f_m$ en el intervalo $[t_{i-1},t_i]$.

Ahora estimemos la oscilación de $f$ en cada intervalo. Si
$x,y\in[t_{i-1},t_i]$, entonces
\[
\begin{aligned}
    |f(x)-f(y)|
    &\leq |f(x)-f_m(x)|
        +|f_m(x)-f_m(y)|
        +|f_m(y)-f(y)|  \\
    &<
        \frac{\varepsilon}{4(b-a)}
        +
        \omega_i(f_m)
        +
        \frac{\varepsilon}{4(b-a)} \\
    &=
        \omega_i(f_m)
        +
        \frac{\varepsilon}{2(b-a)}.
\end{aligned}
\]
Tomando supremo sobre $x,y\in[t_{i-1},t_i]$, obtenemos
\[
    \omega_i(f)
    \leq
    \omega_i(f_m)
    +
    \frac{\varepsilon}{2(b-a)}.
\]
Por tanto,
\[
\begin{aligned}
    \sum_{i=1}^r \omega_i(f)(t_i-t_{i-1})
    &\leq
    \sum_{i=1}^r \omega_i(f_m)(t_i-t_{i-1})
    +
    \frac{\varepsilon}{2(b-a)}
    \sum_{i=1}^r(t_i-t_{i-1}) \\
    &<
    \frac{\varepsilon}{2}
    +
    \frac{\varepsilon}{2(b-a)}(b-a) \\
    &=
    \varepsilon.
\end{aligned}
\]
Por el criterio de integrabilidad mediante oscilaciones, $f$ es integrable.

Ahora probemos la igualdad de integrales. Para $n\geq n_0$,
\[
\begin{aligned}
    \left|
    \int_a^b f(x)\,dx
    -
    \int_a^b f_n(x)\,dx
    \right|
    &=
    \left|
    \int_a^b (f(x)-f_n(x))\,dx
    \right| \\
    &\leq
    \int_a^b |f(x)-f_n(x)|\,dx \\
    &\leq
    \int_a^b \frac{\varepsilon}{4(b-a)}\,dx \\
    &=
    \frac{\varepsilon}{4}.
\end{aligned}
\]
Así,
\[
    \int_a^b f_n(x)\,dx
    \to
    \int_a^b f(x)\,dx.
\]
Luego
\[
    \lim_{n\to\infty}\int_a^b f_n(x)\,dx
    =
    \int_a^b f(x)\,dx.
\]
\end{proof}

\begin{observacion}
Si todas las $f_n$ son continuas, entonces la primera parte de la prueba se
simplifica: por el teorema anterior, el límite uniforme $f$ es continuo; por
tanto, $f$ es integrable.
\end{observacion}

\begin{ejemplo}
La convergencia puntual no basta para pasar el límite dentro de la integral.

Definamos
\[
    f_n:[0,1]\to\R,\qquad f_n(x)=n x^n(1-x^n).
\]
Entonces $f_n(x)\to 0$ puntualmente en $[0,1]$, pero
\[
    \lim_{n\to\infty}\int_0^1 f_n(x)\,dx
    \neq
    \int_0^1 0\,dx.
\]

En efecto, para $0\leq x<1$,
\[
    n x^n\to 0,
\]
y además $f_n(1)=0$. Luego $f_n\to 0$ puntualmente.

Sin embargo,
\[
\begin{aligned}
    \int_0^1 n x^n(1-x^n)\,dx
    &=
    n\int_0^1 x^n\,dx
    -
    n\int_0^1 x^{2n}\,dx \\
    &=
    n\cdot \frac{1}{n+1}
    -
    n\cdot \frac{1}{2n+1} \\
    &=
    \frac{n}{n+1}
    -
    \frac{n}{2n+1}.
\end{aligned}
\]
Por tanto,
\[
    \lim_{n\to\infty}
    \int_0^1 n x^n(1-x^n)\,dx
    =
    1-\frac12
    =
    \frac12.
\]
Pero
\[
    \int_0^1 0\,dx=0.
\]
\end{ejemplo}

%-----------------------------------------------------------
\section{Derivación término a término}
%-----------------------------------------------------------

\begin{grisplano}
\begin{teorema}[Derivación término a término]
Sea $(f_n)$ una sucesión de funciones de clase $C^1$ en $[a,b]$.
Supongamos que:

\begin{enumerate}
    \item existe $c\in[a,b]$ tal que la sucesión numérica
    \[
        (f_n(c))_{n\in\N}
    \]
    converge;

    \item las derivadas $f_n'$ convergen uniformemente en $[a,b]$ a una función
    $g:[a,b]\to\R$.
\end{enumerate}

Entonces $(f_n)$ converge uniformemente en $[a,b]$ a una función $f$ de clase
$C^1$, y
\[
    f'=g.
\]
En particular,
\[
    \left(\lim_{n\to\infty} f_n\right)'
    =
    \lim_{n\to\infty} f_n'.
\]
\end{teorema}
\end{grisplano}

\begin{idea}
No basta pedir que $f_n\to f$ uniformemente para concluir que
$f_n'\to f'$. Lo correcto es controlar las derivadas.

La prueba usa el Teorema Fundamental del Cálculo:
\[
    f_n(x)=f_n(c)+\int_c^x f_n'(t)\,dt.
\]
Luego pasamos al límite. La convergencia de $f_n(c)$ da el valor inicial, y la
convergencia uniforme de $f_n'$ permite pasar el límite dentro de la integral.
\end{idea}

\begin{proof}
Como $(f_n(c))$ converge, existe
\[
    \alpha=\lim_{n\to\infty} f_n(c).
\]
Como $f_n'\to g$ uniformemente y cada $f_n'$ es continua, por el teorema del
límite uniforme de funciones continuas, $g$ es continua.

Para cada $n\in\N$ y cada $x\in[a,b]$, por el Teorema Fundamental del Cálculo,
\[
    f_n(x)
    =
    f_n(c)
    +
    \int_c^x f_n'(t)\,dt.
\]
Como $f_n'\to g$ uniformemente, por el teorema de paso al límite bajo la
integral,
\[
    \lim_{n\to\infty}\int_c^x f_n'(t)\,dt
    =
    \int_c^x g(t)\,dt.
\]
Por tanto, para cada $x\in[a,b]$ existe
\[
    f(x)=\lim_{n\to\infty}f_n(x),
\]
y está dado por
\[
    f(x)
    =
    \alpha
    +
    \int_c^x g(t)\,dt.
\]

Como $g$ es continua, nuevamente por el Teorema Fundamental del Cálculo,
$f$ es derivable y
\[
    f'(x)=g(x)
    \qquad\text{para todo }x\in[a,b].
\]
Además, como $g$ es continua, $f$ es de clase $C^1$.

Falta probar que la convergencia $f_n\to f$ es uniforme. Para todo
$x\in[a,b]$,
\[
\begin{aligned}
    |f_n(x)-f(x)|
    &=
    \left|
    f_n(c)-\alpha
    +
    \int_c^x \bigl(f_n'(t)-g(t)\bigr)\,dt
    \right|\\
    &\leq
    |f_n(c)-\alpha|
    +
    \left|
    \int_c^x \bigl(f_n'(t)-g(t)\bigr)\,dt
    \right|\\
    &\leq
    |f_n(c)-\alpha|
    +
    \int_a^b |f_n'(t)-g(t)|\,dt\\
    &\leq
    |f_n(c)-\alpha|
    +
    (b-a)\sup_{t\in[a,b]}|f_n'(t)-g(t)|.
\end{aligned}
\]
El primer término tiende a $0$ porque $f_n(c)\to\alpha$. El segundo término
tiende a $0$ porque $f_n'\to g$ uniformemente. Por tanto,
\[
    \sup_{x\in[a,b]}|f_n(x)-f(x)|\to 0.
\]
Así, $f_n\to f$ uniformemente.
\end{proof}

\begin{observacion}
Este teorema dice que, bajo sus hipótesis,
\[
    \frac{d}{dx}\left(\lim_{n\to\infty}f_n(x)\right)
    =
    \lim_{n\to\infty}\frac{d}{dx}f_n(x).
\]
Es decir, permite intercambiar límite y derivada.
\end{observacion}

\begin{ejemplo}
Sea
\[
    f_n(x)=\frac{\sen(nx)}{n}.
\]
Entonces
\[
    |f_n(x)|\leq \frac1n
\]
para todo $x\in\R$. Por tanto,
\[
    f_n\to 0
\]
uniformemente en $\R$.

Sin embargo,
\[
    f_n'(x)=\cos(nx).
\]
La sucesión $\cos(nx)$ no converge puntualmente en general, y mucho menos
uniformemente. Por tanto, aunque $f_n\to 0$ uniformemente, no podemos concluir
que
\[
    f_n'\to 0.
\]
\end{ejemplo}
%-----------------------------------------------------------
\section{Series de funciones}
%-----------------------------------------------------------

Sea
\[
    (g_n)_{n\in\N},\qquad g_n:X\to\R,
\]
una sucesión de funciones. La serie de funciones
\[
    \sum_{n=1}^{\infty} g_n
\]
se estudia mediante sus sumas parciales
\[
    s_n(x)=g_1(x)+g_2(x)+\cdots+g_n(x).
\]

\begin{grisplano}
\begin{definicion}[Convergencia puntual y uniforme de una serie]
Decimos que la serie
\[
    \sum_{n=1}^{\infty}g_n
\]
converge puntualmente a $f:X\to\R$ si
\[
    s_n(x)\to f(x)
    \qquad \text{para todo }x\in X.
\]

Decimos que converge uniformemente a $f$ en $X$ si
\[
    s_n\to f
\]
uniformemente en $X$.
\end{definicion}
\end{grisplano}

Si la serie converge a $f$, definimos el resto
\[
    r_n(x)
    =
    \sum_{k=n+1}^{\infty}g_k(x).
\]
Como
\[
    r_n(x)=f(x)-s_n(x),
\]
tenemos la siguiente caracterización.

\begin{grisplano}
\begin{proposicion}[Convergencia uniforme mediante restos]
La serie
\[
    \sum_{n=1}^{\infty}g_n
\]
converge uniformemente a $f$ en $X$ si y solo si
\[
    r_n\to 0
\]
uniformemente en $X$.
\end{proposicion}
\end{grisplano}

\begin{proof}
Como
\[
    r_n=f-s_n,
\]
tenemos
\[
    |r_n(x)|=|f(x)-s_n(x)|.
\]
Por tanto, pedir que $s_n\to f$ uniformemente es exactamente pedir que
$r_n\to 0$ uniformemente.
\end{proof}

%-----------------------------------------------------------
\section{Series de funciones: resultados básicos}
%-----------------------------------------------------------

Los resultados anteriores se aplican a series de funciones usando las sumas
parciales
\[
    s_n=\sum_{k=1}^n f_k.
\]

\begin{grisplano}
\begin{proposicion}[Continuidad de una serie uniforme]
Si
\[
    \sum_{n=1}^\infty f_n
\]
converge uniformemente a $f$ en $X$ y cada $f_n$ es continua en un punto
$a\in X$, entonces $f$ es continua en $a$.
\end{proposicion}
\end{grisplano}

\begin{proof}
Las sumas parciales
\[
    s_n=f_1+\cdots+f_n
\]
son continuas en $a$. Como
\[
    s_n\to f
\]
uniformemente, el límite $f$ es continuo en $a$.
\end{proof}

\begin{grisplano}
\begin{proposicion}[Integración término a término]
Sea
\[
    f_n:[a,b]\to\R
\]
una sucesión de funciones integrables. Si la serie
\[
    \sum_{n=1}^{\infty} f_n
\]
converge uniformemente a $f$ en $[a,b]$, entonces $f$ es integrable y
\[
    \int_a^b f(x)\,dx
    =
    \sum_{n=1}^{\infty}\int_a^b f_n(x)\,dx.
\]
\end{proposicion}
\end{grisplano}

\begin{proof}
Sea
\[
    s_n=\sum_{k=1}^n f_k.
\]
Cada $s_n$ es integrable y
\[
    s_n\to f
\]
uniformemente. Por el teorema de paso al límite bajo el signo integral,
\[
    \int_a^b f(x)\,dx
    =
    \lim_{n\to\infty}\int_a^b s_n(x)\,dx.
\]
Pero
\[
    \int_a^b s_n(x)\,dx
    =
    \int_a^b \sum_{k=1}^n f_k(x)\,dx
    =
    \sum_{k=1}^n\int_a^b f_k(x)\,dx.
\]
Por tanto,
\[
    \int_a^b f(x)\,dx
    =
    \sum_{k=1}^{\infty}\int_a^b f_k(x)\,dx.
\]
\end{proof}

\begin{grisplano}
\begin{proposicion}[Derivación término a término para series]
Sea
\[
    f_n:[a,b]\to\R
\]
una sucesión de funciones de clase $C^1$. Supongamos que:

\begin{enumerate}
    \item la serie numérica
    \[
        \sum_{n=1}^{\infty} f_n(c)
    \]
    converge para algún $c\in[a,b]$;

    \item la serie de derivadas
    \[
        \sum_{n=1}^{\infty} f_n'
    \]
    converge uniformemente en $[a,b]$.
\end{enumerate}

Entonces
\[
    \sum_{n=1}^{\infty} f_n
\]
converge uniformemente a una función $f$ de clase $C^1$, y
\[
    f'
    =
    \sum_{n=1}^{\infty} f_n'.
\]
\end{proposicion}
\end{grisplano}

\begin{proof}
Consideremos las sumas parciales
\[
    s_n=\sum_{k=1}^n f_k.
\]
Entonces
\[
    s_n'=\sum_{k=1}^n f_k'.
\]
Por hipótesis, $(s_n(c))$ converge, porque
\[
    s_n(c)=\sum_{k=1}^n f_k(c),
\]
y la serie $\sum f_k(c)$ converge. Además, $(s_n')$ converge uniformemente,
porque la serie $\sum f_k'$ converge uniformemente.

Aplicando el teorema de derivación término a término a la sucesión $(s_n)$,
concluimos que $(s_n)$ converge uniformemente a una función $f$ de clase $C^1$
y que
\[
    f'
    =
    \lim_{n\to\infty}s_n'
    =
    \sum_{k=1}^{\infty}f_k'.
\]
\end{proof}

%-----------------------------------------------------------
\section{Criterio de Weierstrass}
%-----------------------------------------------------------

\begin{grisplano}
\begin{teorema}[Criterio de Weierstrass]
Sea
\[
    f_n:X\to\R
\]
una sucesión de funciones. Supongamos que existe una serie numérica convergente
\[
    \sum_{n=1}^{\infty} a_n,
    \qquad a_n\geq 0,
\]
tal que
\[
    |f_n(x)|\leq a_n
    \qquad
    \text{para todo }x\in X
    \text{ y todo }n\in\N.
\]
Entonces las series
\[
    \sum_{n=1}^{\infty}|f_n|
    \qquad\text{y}\qquad
    \sum_{n=1}^{\infty}f_n
\]
convergen uniformemente en $X$.
\end{teorema}
\end{grisplano}

\begin{idea}
El criterio compara una serie de funciones con una serie numérica convergente.
Si para todo $x$ los términos de la serie de funciones están dominados por los
términos $a_n$, entonces los restos también quedan dominados:
\[
    \left|\sum_{k=n+1}^{\infty} f_k(x)\right|
    \leq
    \sum_{k=n+1}^{\infty} a_k.
\]
Como el resto numérico puede hacerse menor que $\varepsilon$, el resto funcional
también, uniformemente en $x$.
\end{idea}

\begin{proof}
Para cada $x\in X$, tenemos
\[
    |f_n(x)|\leq a_n.
\]
Como
\[
    \sum_{n=1}^{\infty}a_n
\]
converge, por el criterio de comparación, la serie numérica
\[
    \sum_{n=1}^{\infty}|f_n(x)|
\]
converge para cada $x\in X$. Por tanto, $\sum f_n(x)$ converge absolutamente
para cada $x\in X$.

Ahora probemos la convergencia uniforme. Sea $\varepsilon>0$. Como
\[
    \sum_{n=1}^{\infty}a_n
\]
converge, existe $n_0\in\N$ tal que
\[
    \sum_{k=n+1}^{\infty}a_k<\varepsilon
    \qquad\text{para todo }n\geq n_0.
\]
Definimos el resto
\[
    r_n(x)=\sum_{k=n+1}^{\infty} f_k(x).
\]
Entonces, para todo $x\in X$ y todo $n\geq n_0$,
\[
\begin{aligned}
    |r_n(x)|
    &=
    \left|
    \sum_{k=n+1}^{\infty} f_k(x)
    \right| \\
    &\leq
    \sum_{k=n+1}^{\infty}|f_k(x)| \\
    &\leq
    \sum_{k=n+1}^{\infty}a_k \\
    &<\varepsilon.
\end{aligned}
\]
Como esta estimación no depende de $x$, se tiene
\[
    r_n\to 0
\]
uniformemente en $X$. Luego
\[
    \sum_{n=1}^{\infty} f_n
\]
converge uniformemente.

El mismo argumento aplicado a la serie de funciones no negativas
\[
    \sum_{n=1}^{\infty}|f_n|
\]
demuestra que esta también converge uniformemente.
\end{proof}

\begin{ejemplo}
Consideremos la serie
\[
    \sum_{n=1}^{\infty}\frac{x^n}{n^2}
\]
en el intervalo $[-1,1]$. Para todo $x\in[-1,1]$,
\[
    \left|\frac{x^n}{n^2}\right|
    \leq
    \frac{1}{n^2}.
\]
Como
\[
    \sum_{n=1}^{\infty}\frac{1}{n^2}
\]
converge, por el criterio de Weierstrass la serie
\[
    \sum_{n=1}^{\infty}\frac{x^n}{n^2}
\]
converge uniformemente en $[-1,1]$.
\end{ejemplo}

\begin{observacion}
El criterio de Weierstrass es uno de los instrumentos más importantes para
trabajar con series de funciones. Su ventaja es que reduce un problema de
convergencia uniforme a un problema de convergencia de series numéricas.
\end{observacion}

\end{document}